\section{Experimental Setup}
\label{sec:experiments}

We design experiments to answer three research questions:

\textbf{RQ1:} Does a significant positive correlation exist between TruthfulQA MC1 and AdvGLUE accuracy after controlling for model size?

\textbf{RQ2:} Does model calibration (ECE) correlate with truthfulness and/or robustness?

\textbf{RQ3:} Does calibration moderate the truthfulness-robustness correlation?

\subsection{Models}

We evaluate 14 decoder-only LLMs spanning four families and a range of scales:

\begin{table}[h]
\centering
\begin{tabular}{lll}
\toprule
\textbf{Family} & \textbf{Models} & \textbf{Parameters} \\
\midrule
Pythia & pythia-70m, 160m, 410m, 1b, 1.4b, 2.8b, 6.9b, 12b & 70M--12B \\
Llama-2 & llama-2-7b, 13b, 70b (base) & 7B--70B \\
Mistral & mistral-7b & 7B \\
Falcon & falcon-7b, falcon-40b & 7B--40B \\
\bottomrule
\end{tabular}
\caption{Models evaluated across four families.}
\label{tab:models}
\end{table}

\textbf{Rationale:} Multiple families prevent family-specific artifacts. The size range (70M--70B) enables controlling for scale as a confounder.

\subsection{Benchmarks}

\textbf{TruthfulQA MC1} \cite{lin2022truthfulqa}: 817 questions designed to elicit imitative falsehoods. MC1 format presents one correct answer among distractors; accuracy measures truthful response selection.

\textbf{AdvGLUE} \cite{wang2022advglue}: Adversarial variants of GLUE tasks with word-level perturbations (character swaps, synonym replacements) that preserve semantics. We report average accuracy across subtasks.

\textbf{MMLU} \cite{hendrycks2021mmlu}: 57-subject multiple-choice benchmark used only for ECE computation on a neutral reference task.

\subsection{Evaluation Protocol}

All evaluations use lm-evaluation-harness \cite{gao2023lmeval}:
\begin{itemize}
    \item Batch size: automatically determined per model
    \item Prompt format: default harness templates
    \item Decoding: greedy (temperature=0) for reproducibility
    \item Seed: 42
\end{itemize}

\subsection{Metrics}

\textbf{Primary Metric (RQ1):} Partial Pearson correlation between TruthfulQA MC1 and AdvGLUE accuracy, controlling for $\log(\text{parameters})$.

\textbf{Mechanism Metrics (RQ2-3):}
\begin{itemize}
    \item Expected Calibration Error (ECE): 15-bin ECE on MMLU predictions
    \item ECE-metric correlations: Pearson $r$ between ECE and each trust metric
    \item Moderation test: Fisher's z-test comparing truthfulness-robustness correlation across ECE tertiles
\end{itemize}

\textbf{Statistical Thresholds:}
\begin{itemize}
    \item Significance: $p < 0.05$
    \item Minimum effect size: $|r| > 0.3$
    \item Confidence intervals: 95\% via bootstrap (1000 iterations)
\end{itemize}

\subsection{Hypotheses and Gates}

\begin{table}[h]
\centering
\begin{tabular}{llll}
\toprule
\textbf{ID} & \textbf{Question} & \textbf{Success Criterion} & \textbf{Gate} \\
\midrule
h-e1 & Correlation exists? & $r > 0.3$, $p < 0.05$, CI $> 0$ & MUST\_WORK \\
h-m1 & ECE predicts metrics? & $r < -0.2$ for both & SHOULD\_WORK \\
h-m2 & ECE moderates correlation? & Low-ECE $r >$ High-ECE $r$ & SHOULD\_WORK \\
h-c1 & Pattern holds across model types? & $r > 0.2$ for base and IT & SHOULD\_WORK \\
\bottomrule
\end{tabular}
\caption{Hypotheses and verification gates.}
\label{tab:gates}
\end{table}

The MUST\_WORK gate (h-e1) determines whether the core finding stands; SHOULD\_WORK gates test the calibration mechanism.
